\usepackage[T1]{fontenc}
\usepackage{libertinus}
\usepackage{microtype}
\usepackage{mathtools,amssymb,amsthm}
% Inference rules of the type-theoretic track (Section 1.11 and its boxes).
\usepackage{mathpartir}
\usepackage[nointegrals]{wasysym}
\usepackage{booktabs,tabularx,array,multirow}
\usepackage{xltabular}
\usepackage{enumitem}
\usepackage[dvipsnames,svgnames]{xcolor}
\usepackage{tikz}
\usetikzlibrary{arrows.meta,calc,decorations.pathreplacing,fit,matrix,positioning,shapes.geometric,patterns,backgrounds}

% Chess diagrams (2026-09-04, CHESS-RUNNING-EXAMPLE-PLAN.md phase 1).  The
% pieces are real Staunton figurines from the TeX chess fonts rather than
% letters in discs: vector, already installed, and with no external asset or
% attribution obligation in the build.  Dark squares are hatched by the board
% font, which is the international print convention and is grayscale-safe by
% construction; the book's flat palette is used for everything drawn *over* a
% board, so a chess diagram reads as a chess diagram and an annotation reads as
% one of ours.  Boards are placed inside tikzpictures as nodes, at natural
% size, so the plate discipline of visual-calculus.tex still applies.
\usepackage{chessboard}
\setchessboard{showmover=false,boardfontsize=12pt,labelfontsize=5pt}
\usepackage{adjustbox}
% placeins without the [section] option (2026-09-19, whitespace audit at the
% author's request): the automatic \FloatBarrier at every \section ended a
% page early whenever a figure was still pending, and left forty-odd pages
% with a blank tail of a third or more.  A pending figure now goes to the
% top of the next page and the text fills the page it left; the explicit
% \FloatBarrier commands in the chapters, placed where a plate must precede
% what explains it, are unchanged.
\usepackage{placeins}
\usepackage{etoolbox}
\usepackage{needspace}
% Every \Needspace reserve is capped at four lines (2026-09-19 whitespace
% audit at the author's request).  The chapters ask for up to 26 lines before
% a subsection so that a plate would open with it; that reserved half a page
% and printed as a blank tail whenever less remained.  Four lines keep a
% heading with the start of its text, which the penalty below also does.
% (memoir emulates the needspace package; these are memoir's internals.)
\makeatletter
\newdimen\Art@needdim
\let\Art@needsp@\M@needsp@
\let\Art@sneedsp@\M@sneedsp@
\renewcommand{\M@needsp@}[1]{\Art@needdim=#1\relax
  \ifdim\Art@needdim>4\baselineskip \Art@needdim=4\baselineskip\fi
  \Art@needsp@{\Art@needdim}}
\renewcommand{\M@sneedsp@}[1]{\Art@needdim=#1\relax
  \ifdim\Art@needdim>4\baselineskip \Art@needdim=4\baselineskip\fi
  \Art@sneedsp@{\Art@needdim}}
\makeatother
\usepackage[most]{tcolorbox}
\usepackage{imakeidx}
\usepackage{hyperref}
\usepackage{bookmark}
\usepackage[nameinlink,capitalise,noabbrev]{cleveref}

\setstocksize{9.5in}{6.5in}
\settrimmedsize{\stockheight}{\stockwidth}{*}
\settypeblocksize{6.80in}{4.55in}{*}
\setlrmarginsandblock{0.80in}{1.15in}{*}
\setulmarginsandblock{1.05in}{1.65in}{*}
\checkandfixthelayout
\setlength{\parindent}{1.1em}
% Rubber in \parskip is what lets \flushbottom distribute slack across a page
% instead of dumping it all at the foot.
\setlength{\parskip}{0.15em plus 0.45em minus 0.05em}
\setlength{\emergencystretch}{1.5em}
\linespread{1.17}
\flushbottom
% Flush bottom, but with a bounded reservoir of stretch at the foot.  With
% \flushbottom alone, a page holding a heading, three lines and a tall plate
% has almost no glue to stretch, so the whole deficit lands in one place and
% prints as a hole in the middle of the column.  \raggedbottom cures that and
% costs three times as many short pages.  Giving \@textbottom a finite plus
% makes the foot compete with the in-page glue in proportion: small deficits
% still spread invisibly, large ones drain mostly to the foot where a reader
% expects them.  Measured in tools/measure_whitespace.py.
\makeatletter
\def\@textbottom{\vskip\z@\@plus 10pt}
\makeatother
\widowpenalty=500
\clubpenalty=500
\displaywidowpenalty=500
% This is a continuous teaching edition: do not manufacture blank versos at
% front/main/back-matter transitions. Chapter openings still begin on a new
% page through memoir's openany policy.
\let\cleardoublepage\clearpage
\setlength{\textfloatsep}{12pt plus 3pt minus 4pt}
\setlength{\intextsep}{11pt plus 3pt minus 3pt}
\setlength{\floatsep}{10pt plus 3pt minus 3pt}
\renewcommand{\topfraction}{0.88}
\renewcommand{\bottomfraction}{0.75}
\renewcommand{\textfraction}{0.10}
% 0.76 since the 2026-09-19 whitespace audit (was 0.72): four plates measure
% 72-75% of the text height and were being set alone on a page with a third
% of it blank below; at 0.76 they go to the top of a page and text follows.
% 0.80 was measured too: it produced one overfull page (printed page 392).
\renewcommand{\floatpagefraction}{0.76}
% On a page holding only floats, set them at the top and leave the slack at
% the foot, where a reader expects it (2026-09-19 whitespace audit).  The
% class default centres them, which prints as a hole above the figure.
\makeatletter
\setlength{\@fptop}{0pt}
\setlength{\@dblfptop}{0pt}
\makeatother

\definecolor{ArtNavy}{HTML}{17324D}
\definecolor{ArtBlue}{HTML}{2E6F9E}
\definecolor{ArtTeal}{HTML}{2B8A7E}
\definecolor{ArtOrange}{HTML}{D47A28}
\definecolor{ArtGold}{HTML}{C49A28}
\definecolor{ArtPurple}{HTML}{75539A}
\definecolor{ArtRed}{HTML}{B64A4A}
\definecolor{ArtGreen}{HTML}{4D8452}
\definecolor{ArtPaper}{HTML}{F7F4EC}
\definecolor{ArtGray}{HTML}{5B6570}
\definecolor{ArtLightBlue}{HTML}{EAF3F8}
\definecolor{ArtLightGold}{HTML}{FAF3D7}
\definecolor{ArtLightGreen}{HTML}{EAF4EB}
\definecolor{ArtLightPurple}{HTML}{F1ECF6}
\definecolor{ArtLightRed}{HTML}{FAECEA}

% Shared palette for the recurrence/division tapestry and wheel family.
% Figures 1.24 and 1.25 are the approved rendering baselines.
\definecolor{solarPaper}{HTML}{FCF9F1}
\definecolor{solarInk}{HTML}{283749}
\definecolor{solarMuted}{HTML}{69747D}
\definecolor{solarLine}{HTML}{C8BC9F}
\definecolor{solarGold}{HTML}{DAB453}
\definecolor{solarDeepGold}{HTML}{9D7429}
\definecolor{solarSelect}{HTML}{AC4D43}
\definecolor{solarA}{HTML}{FFD700}
\definecolor{solarB}{HTML}{FFF9E9}
\definecolor{solarC}{HTML}{536477}
\definecolor{solarD}{HTML}{DEA46C}
\definecolor{solarE}{HTML}{81B4A7}
\definecolor{solarF}{HTML}{B1A0CC}
\definecolor{solarG}{HTML}{D78B7D}
\definecolor{solarH}{HTML}{99B18A}
\definecolor{solarI}{HTML}{E6CA7D}
\definecolor{solarJ}{HTML}{95B8CF}

% Glass-bead rendering (2026-09-20, the Glass Bead Game of section 1.6 and
% its plates).  A bead renders one cell or one event; the wire an order or a
% row; a spacer a seam.  Eight phase beads for Z/8 form a hue wheel so that
% opposite residues k and k+4 are complementary colours.  FIGURE-STYLE.md,
% section "Glass-bead family", states the meaning contract.
\definecolor{beadWire}{HTML}{7A6A4F}
\definecolor{beadRim}{HTML}{4E5A66}
\definecolor{beadPost}{HTML}{3E3A33}
\definecolor{ph0}{HTML}{E06C6C}
\definecolor{ph1}{HTML}{E9A15C}
\definecolor{ph2}{HTML}{E3CF62}
\definecolor{ph3}{HTML}{8FC26B}
\definecolor{ph4}{HTML}{5FBFB4}
\definecolor{ph5}{HTML}{6A93D6}
\definecolor{ph6}{HTML}{9B7FD1}
\definecolor{ph7}{HTML}{D57DB5}


\hypersetup{
  bookmarksnumbered=true,
  bookmarksopen=true,
  bookmarksopenlevel=1,
  bookmarksdepth=2,
  pdfpagemode=UseOutlines,
  colorlinks=true,
  linkcolor=ArtNavy,
  citecolor=ArtPurple,
  urlcolor=ArtBlue,
  pdftitle={The Artform Calculus: From Typed Constructions to Quantum Gauge Fields},
  pdfauthor={Sergey Krivov}
}

% Reuse each heading's existing destination: outline-only entries must not
% add printed contents lines, extra anchors, or change the page layout.
\makeatletter
\newcommand{\bookoutline}[2]{%
  \bookmark[level=#1,dest=\@currentHref]{#2}}
% Memoir's numbered-section counter puts its target before the heading
% paragraph. If the heading moves to the next page, that target can stay
% behind. Keep the counter data, but emit its target inside the heading's
% number box, as memhfixc already does for unnumbered section headings.
\newcommand{\book@headingstep}[1]{%
  \let\book@savedraisedlink\Hy@raisedlink
  \let\Hy@raisedlink\@gobble
  \refstepcounter{#1}%
  \let\Hy@raisedlink\book@savedraisedlink}
\patchcmd{\H@old@sectm@m}{\refstepcounter{#1}}{\book@headingstep{#1}}{}
  {\PackageError{artbook}{Cannot patch numbered-section counter}{Check the memoir/hyperref adapter.}}
\patchcmd{\H@old@sectm@m}
  {\protected@edef\@svsec{\@seccntformat{#1}\relax}}
  {\protected@edef\@svsec{\@seccntformat{#1}\noexpand\Hy@SectionAnchorHref{\@currentHref}\relax}}{}
  {\PackageError{artbook}{Cannot attach target to section heading}{Check the memoir/hyperref adapter.}}
\makeatother
\appto\cfttocbeforelisthook{\bookoutline{0}{Contents}}
% The categorical rewrite carries the bibliography beyond printed page 999.
% Reserve room for four-digit page numbers without reducing the text size.
\setpnumwidth{2.5em}
\setrmarg{3.5em}
\indexsetup{othercode={\small\raggedcolumns\bookoutline{0}{Index}}}

\chapterstyle{hangnum}
\renewcommand*{\chapnumfont}{\normalfont\fontsize{44}{48}\selectfont\color{ArtTeal}}
\renewcommand*{\chaptitlefont}{\normalfont\Huge\bfseries\color{ArtNavy}}
\setlength{\beforechapskip}{24pt}
\setlength{\midchapskip}{10pt}
\setlength{\afterchapskip}{26pt}
\setsecheadstyle{\Large\bfseries\color{ArtNavy}}
\setsubsecheadstyle{\large\bfseries\color{ArtTeal}}
\setsubsubsecheadstyle{\normalsize\bfseries\color{ArtPurple}}
\setbeforesecskip{-1.2em plus -0.2em minus -0.2em}
\setaftersecskip{0.7em}

\makepagestyle{artbook}
\makeevenhead{artbook}{\small\color{ArtGray}\thepage}{}{\small\color{ArtGray}THE ARTFORM CALCULUS}
\makeoddhead{artbook}{\small\color{ArtGray}THE ARTFORM CALCULUS}{}{\small\color{ArtGray}\thepage}
\makeevenfoot{artbook}{}{}{}
\makeoddfoot{artbook}{}{}{}
\pagestyle{artbook}
\aliaspagestyle{chapter}{plain}

\newtheoremstyle{artplain}{0.7em}{0.7em}{\itshape}{}{\bfseries\color{ArtNavy}\artfbkhead}{.}{0.5em}{}
\theoremstyle{artplain}
\newtheorem{theorem}{Theorem}[chapter]
\newtheorem{proposition}[theorem]{Proposition}
\newtheorem{lemma}[theorem]{Lemma}
\newtheorem{corollary}[theorem]{Corollary}

\newtheoremstyle{artdef}{0.7em}{0.7em}{}{}{\bfseries\color{ArtTeal}\artfbkhead}{.}{0.5em}{}
\theoremstyle{artdef}
\newtheorem{definition}[theorem]{Definition}
\newtheorem{example}[theorem]{Example}
\newtheorem{principle}[theorem]{Principle}

% Not breakable: a Key idea box carries one self-contained thought, and
% tools/validate_book.py enforces it.  Measured in the 2026-09-19 whitespace
% audit: seven pages keep a tail of a quarter to a third because the box
% missed the foot by a line; that is the price of the rule.
\newtcolorbox{keyidea}[1][]{
  enhanced,colback=ArtLightBlue,colframe=ArtBlue,
  boxrule=0.8pt,arc=2pt,left=7pt,right=7pt,top=6pt,bottom=6pt,
  title={Key idea},fonttitle=\bfseries,#1}
\newtcolorbox{picturefirst}[1][]{
  enhanced,breakable,colback=ArtPaper,colframe=ArtTeal,
  boxrule=0.8pt,arc=2pt,left=7pt,right=7pt,top=6pt,bottom=6pt,
  title={Picture first},fonttitle=\bfseries,artfbk=picturefirst,#1}
\newtcolorbox{rolecheck}[1][]{
  enhanced,colback=ArtLightGold,colframe=ArtGold,
  boxrule=0.8pt,arc=2pt,left=7pt,right=7pt,top=6pt,bottom=6pt,
  title={Role check},fonttitle=\bfseries,#1}
% Retained for historical/pilot documents; removed from the textbook on
% 2026-09-19 at the author's request, with the front-matter legend entry.
\newtcolorbox{worked}[1][]{
  enhanced,breakable,colback=ArtLightGreen,colframe=ArtGreen,
  boxrule=0.8pt,arc=2pt,left=7pt,right=7pt,top=6pt,bottom=6pt,
  title={Worked construction},fonttitle=\bfseries,artfbk=worked,#1}
\newtcolorbox{hypothesis}[1][]{
  enhanced,colback=ArtLightRed,colframe=ArtRed,
  boxrule=0.8pt,arc=2pt,left=7pt,right=7pt,top=6pt,bottom=6pt,
  title={Why the hypothesis matters},fonttitle=\bfseries,#1}
% Retained for historical/pilot documents; removed from the textbook on
% 2026-09-19 at the author's request, with the front-matter legend entry.
\newtcolorbox{formalcorner}[1][]{
  enhanced,breakable,colback=ArtLightPurple,colframe=ArtPurple,
  boxrule=0.8pt,arc=2pt,left=7pt,right=7pt,top=6pt,bottom=6pt,
  title={Formal corner},fonttitle=\bfseries,artfbk=formalcorner,#1}
% modelboundary and hypothesis are deliberately NOT breakable.  Making them so
% was tried on 2026-09-04 to close a hole on the page before a long Model
% boundary box, and it drew a Key idea box on top of a Try it box two pages
% earlier: with breakable in play tcolorbox mispositioned a later atomic box
% whose float had moved.  tools/find_overprint.py catches that class of fault;
% it found this one, and reverting breakable cleared it.
% Retained for historical/pilot documents; removed from the textbook on 2026-09-18.
\newtcolorbox{modelboundary}[1][]{
  enhanced,colback=white,colframe=ArtGray,
  frame style={dash pattern=on 3pt off 2pt},
  boxrule=0.8pt,arc=2pt,left=7pt,right=7pt,top=6pt,bottom=6pt,
  title={Model boundary},fonttitle=\bfseries,#1}
% Making `rolecheck' and `checkpoint' breakable was measured on 2026-09-04
% while closing the whitespace gaps: no overprints appeared, but the pages with
% more than an inch of conspicuous white went 21 to 21 with both and to 22 with
% `rolecheck' alone, so neither earns the risk recorded above.  `checkpoint'
% must stay atomic in any case: the frontmatter promises the answer is just
% below the question.
\newtcolorbox{checkpoint}[1][]{
  enhanced,colback=white,colframe=ArtOrange,
  boxrule=0.7pt,arc=2pt,left=7pt,right=7pt,top=5pt,bottom=5pt,
  title={Try it},fonttitle=\bfseries,#1}
% "Why it is built this way" answers a design question: why the definition has
% the shape it has, what the obvious alternative would be, and what goes wrong
% with it. That is a different question from `hypothesis', which shows what
% breaks when a stated condition is dropped. Added 2026-09-04 for chapter 4,
% where the design questions are the ones a good reader actually asks.
\newtcolorbox{design}[1][]{
  enhanced,breakable,colback=ArtLightGreen,colframe=ArtGreen,
  boxrule=0.8pt,arc=2pt,left=7pt,right=7pt,top=6pt,bottom=6pt,
  title={Why it is built this way},fonttitle=\bfseries,artfbk=design,#1}
% The book no longer opens its chapters with this box: the author removed all
% thirteen on 2026-09-15 to shorten the edition, and the front matter's list of
% boxes lost its entry with them.  The definition stays because the pilot
% documents (quantum-pilots/pilots.tex, matter-pilots) input this file and
% still use it.  Nothing in chapters/ does.
\newtcolorbox{chaptergoals}[1][]{
  enhanced,breakable,colback=ArtPaper,colframe=ArtNavy,
  boxrule=1pt,arc=2pt,left=8pt,right=8pt,top=7pt,bottom=7pt,
  title={What you will learn},fonttitle=\bfseries,artfbk=chaptergoals,#1}

% A heading is never left at the foot of a page with nothing under it.
%
% This is done with a PENALTY, not by reserving space, and the distinction was
% expensive to learn.  Both \needspace and \Needspace* reserve room by putting
% glue on the vertical list ahead of the heading, and glue ahead of a heading
% that sits under a tcolorbox goes wrong: the reserve is measured against
% \pagegoal-\pagetotal before the box has finished contributing, so the break
% is not taken and the heading is set over the box's last line.  That was the
% overlap at section 4.5, at the Part B divider, on the "Three reading routes"
% page, and on the first two lines of the table of contents -- intermittent,
% because it depends on where the break falls.  Reserving space also cannot
% help against an unplaced float, whose height is likewise unknown at the time
% the reserve is measured.  Do not reintroduce either macro here.
%
% The penalty route asks the paragraph after a heading to keep its first three
% lines together.  \@afterheading already sets \clubpenalty to \@M for the
% first line; eTeX's \clubpenalties extends that to the second, so the earliest
% legal break is after the third line and a heading always carries at least
% that much with it.  The setting is scoped to the paragraph that follows a
% heading -- \everypar restores the ordinary \clubpenalty at the next one -- so
% it costs nothing anywhere else in the book.  Applying \clubpenalties to every
% paragraph was measured and rejected: it cost five pages.
% (Diagnosed 2026-09-04.)
\makeatletter
\def\@afterheading{%
  \@nobreaktrue
  \everypar{%
    \if@nobreak
      \@nobreakfalse
      \clubpenalties 3 \@M \@M \@clubpenalty
      \if@afterindent \else {\setbox\z@\lastbox}\fi
    \else
      \clubpenalties 1 \@clubpenalty
      \everypar{}%
    \fi}}
\makeatother

\NewDocumentCommand{\bookdivision}{s m m}{%
  % A part divider introduces what follows it, so it must not be the last
  % thing on a page. Reserve the bar plus the heading and a few lines of the
  % section it opens; otherwise it strands at a page foot and leaves the rest
  % of that page empty. (Added 2026-09-04 after chapter 4's Part B divider did
  % exactly that.)
  \par\addvspace{1.25em}
  \IfBooleanTF{#1}{%
  \begin{tcolorbox}[
    enhanced,colback=ArtNavy,colframe=ArtNavy,boxrule=0pt,arc=2pt,
    left=9pt,right=9pt,top=7pt,bottom=7pt,
    after={\csname tcb@parfillskip@check\endcsname\par\nobreak\smallskip}]
    \color{white}\sffamily
    {\Large\bfseries Part #2}\hspace{0.8em}{\large #3}
  \end{tcolorbox}}{%
  \begin{tcolorbox}[
    enhanced,colback=ArtNavy,colframe=ArtNavy,boxrule=0pt,arc=2pt,
    left=9pt,right=9pt,top=7pt,bottom=7pt]
    \color{white}\sffamily
    {\Large\bfseries Part #2}\hspace{0.8em}{\large #3}
  \end{tcolorbox}}%
  \phantomsection
  \addcontentsline{toc}{section}{Part #2: #3}
  \IfBooleanTF{#1}{%
  % The starred form keeps the divider with the heading it introduces
  % (2026-09-23, chapter 5, whose Part B divider stranded at a page foot).
  % Three breaks are closed. The box's after code is tcolorbox's default
  % (\par\smallskip) with a \nobreak before the skip. The nobreak flag makes
  % the \section that follows add neither its break penalty nor its skip, so
  % that skip is laid here: the same \beforesecskip, and the spacing is
  % unchanged. And the heading's \color would put a colour whatsit on the
  % vertical list ahead of the heading paragraph, where the \parskip glue after
  % it is a legal break; for this one heading the paragraph starts first, so
  % the colour lands inside it, and the style restores itself at once. The
  % plain form is unchanged, so no other chapter's pages move.
  \nobreak\vskip-\beforesecskip\relax
  \csname @nobreaktrue\endcsname
  \global\let\bookdivisionsecheadstyle\secheadstyle
  \gdef\secheadstyle{\noindent\bookdivisionsecheadstyle
    \global\let\secheadstyle\bookdivisionsecheadstyle}}{%
  \nobreak\vspace{0.35em}}%
}

\newenvironment{exercises}{
  \section*{Exercises}
  \phantomsection
  \addcontentsline{toc}{section}{Exercises}
  % Shrinkable itemsep: a long exercise list may tighten slightly rather than
  % spill its last one or two items onto an otherwise empty page before the
  % next chapter. 15 items give about four lines of slack.
  \begin{enumerate}[label=\textbf{\thechapter.\arabic*},leftmargin=2.8em,
    itemsep={0.8em plus 0.2em minus 0.32em}]
}{\end{enumerate}}

% Difficulty dots sit flush right on the last line and never wrap alone.
\newcommand{\difficulty}[1]{\unskip\nobreak\hfil\penalty50\hskip1em\hbox{}\nobreak\hfil
  \mbox{\color{ArtOrange}\textbf{#1}}%
  {\parfillskip=0pt\finalhyphendemerits=0\par}}
\newcommand{\role}[1]{\textsf{\bfseries\color{ArtPurple}#1}}
% \term[index entry]{Printed term}: the optional argument is the index entry,
% for a term printed in a shorter or inflected form (2026-09-23); a capitalised
% entry keeps its sort key, as in \term[hilbert space@Hilbert space]{Hilbert space}.
\NewDocumentCommand{\term}{o m}{\textbf{#2}\IfNoValueTF{#1}{\index{#2}}{\index{#1}}}
% \lawforward opens a paragraph that says what a construction of Chapter 2
% becomes in the definition of a physical law in Chapter 5 (author, 2026-09-24).
\newcommand{\lawforward}{\par\noindent{\sffamily\bfseries\color{ArtTeal}In a physical law.}\enspace}
\newcommand{\cells}{\AFCell}
\newcommand{\supp}{\operatorname{Supp}}
\newcommand{\pos}{\AFPos}
\newcommand{\state}{\AFState}
\newcommand{\val}{\operatorname{Val}}
\newcommand{\RowShape}{\operatorname{RowShape}}
\newcommand{\TempSpec}{\operatorname{TempSpec}}
\newcommand{\assemble}{\AFAssemble}
\newcommand{\glue}{\AFGlue}
\newcommand{\subst}{\AFSubstitute}
\newcommand{\Align}{\operatorname{Align}}
\newcommand{\Emb}{\operatorname{Emb}}
\newcommand{\Place}{\operatorname{Place}}
% Positive comparison (chapter 3, 2026-09-05): positive reduct, atom image,
% credited atoms. Plain operators; the witness types are in the registry.
\newcommand{\AtP}{\operatorname{At}^{+}}
\newcommand{\aim}{\operatorname{aim}}
\newcommand{\Pres}{\AFPresentation}
\newcommand{\Aut}{\operatorname{Aut}}
\newcommand{\Stab}{\operatorname{Stab}}
\newcommand{\Fix}{\operatorname{Fix}}
\newcommand{\Cont}{\operatorname{Cont}}
\newcommand{\curl}{\operatorname{curl}}
\newcommand{\divg}{\operatorname{div}}
\newcommand{\Hol}{\operatorname{Hol}}
% Added for the gauge expansion (GAUGE-EXPANSION-PLAN.md phase 0).
% Plain mathematical operators only. New *public constructor* names
% (localization, torsor, link table) belong in constructor-registry.tex
% with a CONSTRUCTOR-NAME-LEDGER.md entry, and are added in the phase
% that first writes the prose using them.
\newcommand{\Hom}{\operatorname{Hom}}
\newcommand{\End}{\operatorname{End}}
\newcommand{\Tr}{\operatorname{Tr}}
\newcommand{\Irr}{\operatorname{Irr}}
\newcommand{\id}{\operatorname{id}}
\newcommand{\Z}{\mathbb Z}
\newcommand{\N}{\mathbb N}
\newcommand{\Q}{\mathbb Q}
\newcommand{\R}{\mathbb R}
\newcommand{\F}{\mathbb F}
\newcommand{\C}{\mathbb C}
\newcommand{\Zmod}[1]{\Z/#1\Z}
\newcommand{\one}{\mathbf 1}
\newcommand{\blank}{\square}
\newcommand{\dd}{\,\mathrm d}
\newcommand{\angles}[1]{\langle #1\rangle}

% Textbook adapters to the corpus-wide CN-1 public vocabulary.
% Definitions and interfaces are owned by ../NOTATION.md and ../CONSTRUCTORS.md.
% Historical source spellings remain in CONSTRUCTOR-NAME-LEDGER.md.

% Sorts and role-indexed families.
\newcommand{\OrderedForm}{\AFOrderedForm}
\newcommand{\TimeForm}{\AFTimeForm}
\newcommand{\PointedTimeForm}{\AFPointedTimeForm}
\newcommand{\TapestryForm}{\AFTapestryForm}
\newcommand{\CellForm}{\AFCellForm}
\newcommand{\ChromoForm}{\AFChromoForm}
\newcommand{\Presentation}{\AFPresentation}
\newcommand{\State}{\AFState}
\newcommand{\Witness}{\AFWitness}
\newcommand{\Script}{\AFScript}
\newcommand{\Scheduler}{\AFScheduler}
\newcommand{\Observable}{\AFObservable}
\newcommand{\Rendering}{\AFRendering}

% Finite kernel and ordered arithmetic.
\newcommand{\FinCode}{\AFFinCode}
\newcommand{\History}{\AFHistory}
\newcommand{\Empty}{\AFEmpty}
\newcommand{\Atom}{\AFAtom}
\newcommand{\TypedAtom}{\AFTypedAtom}
\newcommand{\ValuedAtom}{\AFValuedAtom}
\newcommand{\Index}{\AFIndex}
\newcommand{\Replicate}{\AFReplicate}
\newcommand{\Assemble}{\AFAssemble}
\newcommand{\OrderedAssemble}{\AFAssemble^{\mathrm{ord}}}
\newcommand{\Glue}{\AFGlue}
\newcommand{\Substitute}{\AFSubstitute}
\newcommand{\Normalize}{\AFNormalize}
\newcommand{\NatForm}{\AFNatForm}
\newcommand{\Rod}{\AFRod}
\newcommand{\SegmentedRod}{\AFSegmentedRod}
\newcommand{\ProductRod}{\AFProductRod}
\newcommand{\GridView}{\AFGridView}
\newcommand{\Wheel}{\AFWheel}
\newcommand{\RatioForm}{\AFRatioForm}
\newcommand{\PerfectAlign}{\AFPerfectAlign}
\newcommand{\AddWitness}{\AFAddWitness}
\newcommand{\MulWitness}{\AFMulWitness}

% Time, tapestry, and dynamics.
\newcommand{\LinearTime}{\AFLinearTime}
\newcommand{\CycleTime}{\AFCycleTime}
\newcommand{\PointTime}{\AFPointTime}
\newcommand{\Clock}{\AFClock}
\newcommand{\SampleClock}{\AFSampleClock}
\newcommand{\Orb}{\AFOrb}
\newcommand{\RectTapestry}{\AFRectTapestry}
\newcommand{\EventTapestry}{\AFEventTapestry}
\newcommand{\TemporalCode}{\AFTemporalCode}
\newcommand{\ProcessProfile}{\AFProcessProfile}
\newcommand{\SchedulerPolicy}{\AFSchedulerPolicy}
\newcommand{\CompatibleBatch}{\AFCompatibleBatch}
\newcommand{\Sweep}{\AFSweep}
\newcommand{\Trajectory}{\AFTrajectory}
\newcommand{\TraceTapestry}{\AFTraceTapestry}
\newcommand{\ConcurrencyWitness}{\AFConcurrencyWitness}

% Cellular and pattern geometry.
\newcommand{\Canvas}{\AFCanvas}
\newcommand{\SquareGrid}{\AFSquareGrid}
\newcommand{\Quotient}{\AFQuotient}
\newcommand{\AttachCell}{\AFAttachCell}
\newcommand{\Refine}{\AFRefine}
% Recognition modes. The ledger keeps all three: Match is the general family
% of profile correspondences, Occurrence its strong mode, Placement the
% preserve-only mode. They share a register because they share a family.
\newcommand{\Match}{\AFMatch}
\newcommand{\GeometricForm}{\AFGeometricForm}
\newcommand{\RecognitionPattern}{\AFRecognitionPattern}
\newcommand{\Occurrence}{\AFOccurrence}
\newcommand{\Placement}{\AFPlacement}
% Positive comparison and observation. Added 2026-09-05 with the
% science-of-art foundation (SCIENCE-OF-ART-FOUNDATION-GRAFT.md); the naming
% decisions are recorded in CONSTRUCTOR-NAME-LEDGER.md.
\newcommand{\Instantiation}{\AFInstantiation}
\newcommand{\Observe}{\AFObserve}
\newcommand{\SampledPicture}{\AFSampledPicture}
\newcommand{\ObservationProtocol}{\AFProtocol}
\newcommand{\FeatureState}{\AFFeatureState}
\newcommand{\Token}{\AFToken}
\newcommand{\PatchWitness}{\AFPatchWitness}
\newcommand{\Realization}{\AFRealization}
\newcommand{\Realize}{\AFRealize}
\newcommand{\Run}{\AFRun}
\newcommand{\Realizer}{\AFRealizer}
\newcommand{\Interpret}{\AFInterpret}
% Interpretation witnesses.  The generic core is a dependent pair of a
% realizer and its profile-specific law certificate; the named forms below
% are conservative specializations used by the textbook.
\newcommand{\IntWit}{\AFIntWit}
\newcommand{\Sem}{\AFSem}
\newcommand{\PointedInt}{\AFPointedInt}
\newcommand{\ObsIntWit}{\AFObsIntWit}
\newcommand{\MatchIntWit}{\AFMatchIntWit}
\newcommand{\CoordIntWit}{\AFCoordIntWit}
\newcommand{\GaugeObsIntWit}{\AFGaugeObsIntWit}
\newcommand{\GaugeSafeRecIntWit}{\AFGaugeSafeRecIntWit}
\newcommand{\SubstitutionStage}{\AFSubstitutionStage}

% Gauge constructors. Added 2026-09-04 with the general gauge sections of
% chapter 4; the naming decision is recorded in CONSTRUCTOR-NAME-LEDGER.md 4.8.
\newcommand{\Torsor}{\AFTorsor}
\newcommand{\PoseOrbit}{\AFPoseOrbit}
\newcommand{\Localize}{\AFLocalize}
\newcommand{\LinkTable}{\AFLinkTable}

% Measure and physical application layer.
\newcommand{\Tally}{\AFTally}
\newcommand{\Content}{\AFContent}
\newcommand{\CoefficientProfile}{\AFCoefficientProfile}
\newcommand{\Field}{\AFField}
\newcommand{\Connection}{\AFConnection}
\newcommand{\Chromo}{\AFChromo}
\newcommand{\ChromoState}{\AFChromoState}
\newcommand{\RewriteApplication}{\AFRewriteApplication}
\newcommand{\LawSpec}{\AFLawSpec}
\newcommand{\MaxwellState}{\AFMaxwellState}

% Judgments. These macros make the role visible in display mathematics.
\newcommand{\forms}[1]{\vdash #1\;\mathsf{form}}
\newcommand{\states}[2]{#1:#2}
\newcommand{\witnesses}[3]{#1:#2\rightsquigarrow #3}
\newcommand{\scripts}[3]{#1:#2\longrightarrow #3}
\newcommand{\constructs}[3]{#1\mathrel{\xrightarrow{\ #2\ }}#3}
\newcommand{\reduces}[3]{#1\mathrel{\xrightarrow[\mathrm{forget}]{\ #2\ }}#3}

% Temporal counting and conservative record interfaces.
\newcommand{\Cyclarchy}{\AFCyclarchy}
\newcommand{\OccurrenceHistory}{\AFOccurrenceHistory}
\newcommand{\TickRealization}{\AFTickRealization}
\newcommand{\DurationRecord}{\AFDurationRecord}
\newcommand{\EventAssociation}{\AFEventAssociation}
\newcommand{\CalibrationRecord}{\AFCalibrationRecord}
\newcommand{\Elapsed}{\AFElapsed}

% Qualified spatial measurement interfaces.
\newcommand{\MeasurandSpec}{\AFMeasurandSpec}
\newcommand{\QuantityValue}{\AFQuantityValue}
\newcommand{\MetricContentRecord}{\AFMetricContentRecord}
\newcommand{\SpatialAssociation}{\AFSpatialAssociation}
\newcommand{\SpatialCalibrationRecord}{\AFSpatialCalibrationRecord}
\newcommand{\SpatialMeasurementWitness}{\AFSpatialMeasurementWitness}
\newcommand{\SpatialRead}{\AFSpatialRead}


\tikzset{
  book diagram/.style={
    font=\small,line cap=round,line join=round,
    show background rectangle,inner frame sep=4pt,
    background rectangle/.style={
      fill=ArtPaper!72,draw=ArtGray!24,line width=.45pt,rounded corners=4pt}},
  cell/.style={draw=ArtNavy,thick,minimum width=7mm,minimum height=7mm,inner sep=1pt},
  witness/.style={draw=ArtGold,very thick,-{Latex[length=2mm]}},
  transform/.style={draw=ArtPurple,very thick,-{Latex[length=2.2mm]}},
  incidence/.style={draw=ArtNavy,thick},
  control/.style={fill=ArtLightBlue,draw=ArtBlue,thick},
  state/.style={fill=ArtLightGreen,draw=ArtGreen,thick},
  derived/.style={fill=ArtLightPurple,draw=ArtPurple,thick},
  mirror/.style={draw=ArtRed,very thick},
  root/.style={draw=ArtBlue,very thick,-{Latex[length=2mm]}},
  lightgrid/.style={draw=ArtGray!35,thin}
}

% Keep the visual language local to actual figures; tcolorbox and theorem
% internals also use TikZ and must not inherit the diagram backdrop.
\AtBeginEnvironment{figure}{%
  \tikzset{every picture/.append style={book diagram}}%
}

% Reusable semantic components for the textbook's visual calculus.
% The styles encode mathematical roles and must remain distinguishable in
% grayscale.  Figures may specialize geometry but not the role grammar.
%
% Plate discipline (2026-09-02 repair): plates are drawn at natural size and
% must fit the 4.55in text block; no \resizebox or adjustbox max-width
% shrinking is permitted around a plate, so node text never prints below the
% \scriptsize used here.  Arrow labels always carry the `arrow label' backing
% so they never overprint arrows, frames, or neighbouring nodes.

\tikzset{
  constructor arrow/.style={draw=ArtTeal,ultra thick,-{Latex[length=2.6mm]},
    double=white,double distance=.7pt},
  witness arrow/.style={draw=ArtGold!70!black,very thick,dashed,
    -{Latex[length=2.2mm]}},
  script arrow/.style={draw=ArtPurple,very thick,-{Latex[length=2.5mm]}},
  reduct arrow/.style={draw=ArtGray,thick,densely dotted,
    -{Latex[length=2.1mm]}},
  dependency arrow/.style={draw=ArtGray!80,thin,-{Latex[length=1.7mm]}},
  % A dependency the role policy forbids outright (the Ren row of the policy
  % matrix).  Dotted red plus a crossed marker so it stays legible in grayscale
  % and can never be mistaken for the merely gate-mediated `witness arrow'.
  forbidden arrow/.style={draw=ArtRed,thick,densely dotted,
    -{Latex[length=2.1mm]}},
  forbidden mark/.style={circle,draw=ArtRed,fill=white,line width=.5pt,
    inner sep=.6pt,font=\tiny\bfseries,text=ArtRed},
  arrow label/.style={fill=ArtPaper!72,inner sep=1.5pt,font=\scriptsize,
    align=center},
  plain label/.style={fill=white,inner sep=1.5pt,font=\scriptsize,
    align=center},
  type node/.style={draw=ArtNavy,rounded corners=2pt,fill=white,
    minimum height=7mm,align=center,font=\scriptsize,inner sep=3pt},
  guard node/.style={draw=ArtGold!70!black,diamond,aspect=2.0,
    fill=ArtLightGold,align=center,font=\scriptsize,inner sep=1.5pt},
  process lane/.style={draw=ArtGray!55,rounded corners=2pt,fill=white,
    minimum height=6mm,align=left,font=\scriptsize},
  formal cell/.style={cell,font=\scriptsize,fill=white},
  wide cell/.style={cell,font=\scriptsize,fill=white,minimum width=12mm},
  provenance left/.style={formal cell,fill=ArtLightBlue},
  provenance right/.style={formal cell,fill=ArtLightGold},
  retained overlay/.style={draw=ArtOrange,very thick},
  rendering boundary/.style={draw=ArtGray,dashed,rounded corners=3pt,
    fill=ArtPaper!55,align=center,font=\scriptsize},
  seam/.style={draw=ArtRed,very thick}
}

\newcommand{\rolebadge}[1]{%
  \tikz[baseline=(r.base)]\node[draw=ArtPurple,rounded corners=2pt,
  fill=ArtLightPurple,inner xsep=4pt,inner ysep=1.4pt,font=\scriptsize\bfseries]
  (r){#1};}

\newcommand{\rolestrip}[1]{%
  \par\nobreak\smallskip\nobreak\noindent
  \colorbox{ArtPaper}{\parbox{\dimexpr\linewidth-2\fboxsep\relax}{\scriptsize
  \raggedright\textbf{ROLE}\enspace #1\par}}\par\smallskip}

% Formation and carrier lines are typeset at \small without any shrinking.
% Long judgments break across lines with `\\' and lead every row with `&'.
% The label sits on its own line so the judgment gets the full measure, and
% the mathematics is set one size down: a formation line must never overflow.
\newcommand{\formationline}[1]{%
  \par\smallskip\noindent
  \colorbox{ArtLightBlue}{\parbox{\dimexpr\linewidth-2\fboxsep\relax}{%
  \raggedright\small\textbf{Formation.}\par\nobreak\vspace{1pt}\nobreak
  \footnotesize$\displaystyle\begin{aligned}[t]#1\end{aligned}$\par}}\par\smallskip}

\newcommand{\carrierline}[1]{%
  \par\smallskip\noindent
  \colorbox{ArtLightGreen}{\parbox{\dimexpr\linewidth-2\fboxsep\relax}{%
  \raggedright\small\textbf{Carrier.}\enspace
  \footnotesize$\displaystyle\begin{aligned}[t]#1\end{aligned}$\par}}\par\smallskip}

\newcommand{\forgetline}[1]{%
  \par\smallskip\noindent
  \colorbox{ArtPaper}{\parbox{\dimexpr\linewidth-2\fboxsep\relax}{\small
  \textbf{Forgetting boundary.}\enspace #1}}\par\smallskip}

% Study headings: unnumbered, one size below \subsection, kept with the
% text that follows them (no page break after the title or its rule).
\makeatletter
\newcommand{\studyheading}[2]{%
  % No \needspace (2026-09-29): a reserve measured after the role strip of the
  % study above misjudged the room and left printed page 347 a third blank;
  % \@afterheading keeps the heading with three lines of its text (style.tex).
  \par\addvspace{1.1em}\noindent
  {\normalsize\bfseries\color{ArtNavy}#2}\par\nobreak\smallskip\nobreak
  \noindent\textcolor{#1}{\rule{\linewidth}{1.1pt}}\par\nobreak\smallskip
  \@afterheading}
\makeatother

\newenvironment{constructorstudy}[1]{\studyheading{ArtTeal}{#1}}
  {\par\addvspace{.8em}}
\newenvironment{witnessstudy}[1]{\studyheading{ArtGold!75!black}{#1}}
  {\par\addvspace{.8em}}
\newenvironment{processstudy}[1]{\studyheading{ArtPurple}{#1}}
  {\par\addvspace{.8em}}

% A float-safe worked derivation.  The general ``worked'' box is breakable,
% hence it cannot contain figure floats.  Visual derivations deliberately use
% this open rule treatment so their numbered plates remain ordinary floats.
\newenvironment{guidedderivation}[1]{\studyheading{ArtGreen}{#1}}
  {\par\addvspace{.8em}}

\newcommand{\schemainhabitantcaption}[2]{%
  \par\smallskip
  \begin{center}\scriptsize
  \begin{tabularx}{.96\linewidth}{>{\bfseries\color{ArtTeal}}p{.18\linewidth}X}
  schema & #1\\
  inhabitant & #2
  \end{tabularx}
  \end{center}}

\newcommand{\visualaudit}[4]{%
  \begin{center}\scriptsize
  \begin{tabularx}{.96\linewidth}{>{\bfseries}p{.19\linewidth}X}
  creates & #1\\
  retains & #2\\
  effect & #3\\
  does not create & #4
  \end{tabularx}
  \end{center}}

% ------------------------------------------------- a second chess piece set
% fig:chess-renderings needs two renderings of one position that a reader can
% tell apart at a glance.  Panel (a) is the convention of printed chess
% literature: Staunton figurines on a hatched board, set from the TeX chess
% font by the chessboard package.  No second board font shipped with those
% packages differs from the first by enough to make the point -- skak and
% skaknew are near-identical at plate size, and berlin silently falls back to
% skak -- and no SVG converter is installed to bring in an outside set, which
% the plate contract would in any case refuse as non-editable.  So the second
% set is drawn here: flat squares in the book's palette, geometric pieces.
% White pieces are outlined and black pieces solid, so tone and not hue
% carries the side.  Only the four kinds the running position uses are
% defined; a fifth would be dead code until a plate needs it.
\newcommand{\flatboard}[1]{%
  \foreach \bdf in {0,...,7}{\foreach \bdr in {0,...,7}{%
    \ifodd\numexpr\bdf+\bdr\relax
      \fill[white] (\bdf*#1,\bdr*#1) rectangle (\bdf*#1+#1,\bdr*#1+#1);
    \else
      \fill[ArtGray!30] (\bdf*#1,\bdr*#1) rectangle (\bdf*#1+#1,\bdr*#1+#1);
    \fi}}%
  \draw[ArtNavy,thick] (0,0) rectangle (8*#1,8*#1);}

\newcommand{\flatlabels}[1]{%
  \foreach \bdf/\bdl in {0/a,1/b,2/c,3/d,4/e,5/f,6/g,7/h}
    \node[font=\tiny,ArtGray,anchor=north,yshift=-.6mm]
      at (\bdf*#1+0.5*#1,0) {\vphantom{b}\bdl};
  \foreach \bdr in {0,...,7}{%
    \pgfmathtruncatemacro{\bdn}{\bdr+1}%
    \node[font=\tiny,ArtGray,anchor=base east]
      at (-0.22*#1,\bdr*#1+0.5*#1-0.055*#1) {\bdn};}}

% \legendswatch{<colour>} -- a framed 2.6 mm colour square for a legend row,
% set inline before its word so the word can never run into the next swatch
% (fig:zodiac-wheel).  Lowered a little so it sits on the x-height.
\newcommand{\legendswatch}[1]{%
  \raisebox{-0.45mm}{\setlength{\fboxsep}{0pt}\setlength{\fboxrule}{0.4pt}%
    \fcolorbox{ArtNavy}{#1}{\rule{0pt}{2.6mm}\rule{2.6mm}{0pt}}}}

% The two hand-drawn piece sets that used to live here -- \flatpiece, solid
% bodies, and \strokepiece, stems with finials -- were removed on 2026-09-04.
% fig:chess-renderings now draws both of its panels with the `chessboard'
% package: figurines on the left, \wdisc/\bdisc lettered discs on the right.
% Drawing chessmen by hand is not worth it at plate size; the font is better
% at it, and two drawing systems in one plate never looked right.
% \flatboardplain{<square side>} -- the same board with the colouring removed,
% for the panel of fig:board-colour-group that shows what a profile sees when
% it reads incidence only.  Gridlines, because without them there is nothing
% to see; they render the incidence and are not carried data.
\newcommand{\flatboardplain}[1]{%
  \fill[white] (0,0) rectangle (8*#1,8*#1);
  \foreach \bdk in {1,...,7}{%
    \draw[ArtGray!45,thin] (\bdk*#1,0)--(\bdk*#1,8*#1) (0,\bdk*#1)--(8*#1,\bdk*#1);}
  \draw[ArtNavy,thick] (0,0) rectangle (8*#1,8*#1);}

% \flatboardbw{<square side>} -- the plain board of \flatboardplain with the
% palette applied and painted, black against white rather than the printer's
% gray.  Used for the right-hand board of fig:board-construction, so that both
% boards in that plate are the same drawing and only the marking differs.
% No gridlines: the fills already show every edge.
\newcommand{\flatboardbw}[1]{%
  \foreach \bdf in {0,...,7}{\foreach \bdr in {0,...,7}{%
    \ifodd\numexpr\bdf+\bdr\relax
      \fill[white] (\bdf*#1,\bdr*#1) rectangle (\bdf*#1+#1,\bdr*#1+#1);
    \else
      \fill[black] (\bdf*#1,\bdr*#1) rectangle (\bdf*#1+#1,\bdr*#1+#1);
    \fi}}%
  \draw[ArtNavy,thick] (0,0) rectangle (8*#1,8*#1);}

% Letters in discs, for the second panel of fig:chess-renderings.  They are
% placed on a standard `chessboard' board through pgfstyle=text, so the board
% itself -- squares, hatching, border, coordinates -- is the package's own and
% identical to the panel beside it.  `show background rectangle=false' is
% required: style.tex gives every picture the book's paper backing, which on a
% disc this small would print as a grey card behind each piece.
\newcommand{\wdisc}[1]{\tikz[baseline=-0.62ex,show background rectangle=false]
  \node[circle,draw=ArtNavy,fill=white,line width=.55pt,inner sep=0pt,
        minimum size=3.5mm,font=\bfseries\scriptsize,text=ArtNavy]{#1};}
\newcommand{\bdisc}[1]{\tikz[baseline=-0.62ex,show background rectangle=false]
  \node[circle,draw=ArtNavy,fill=ArtNavy,line width=.55pt,inner sep=0pt,
        minimum size=3.5mm,font=\bfseries\scriptsize,text=white]{#1};}


% ------------------------------------------------- the glass-bead rendering
% Added 2026-09-20 for the Glass Bead Game (section 1.6).  The styles are
% prefixed gb so that no figure-local style of the same name is overridden.
% A bead renders one cell or one event and its colour a label or value; a wire
% renders a linear order or one row of a tapestry, ended by knots or posts; a
% spacer renders a seam, thick for the last constructor's seam and thin for an
% inherited one; a long bead renders one value held over several positions.
% Nothing drawn by these macros is data: no rule reads a bead.
\tikzset{
  gbwire/.style={draw=beadWire,line width=1.1pt,line cap=round},
  gbwire faint/.style={draw=beadWire!55,line width=.8pt,line cap=round},
  gbwire arrow/.style={draw=beadWire,line width=1.1pt,line cap=round,
    -{Latex[length=1.7mm]}},
  gbpost/.style={draw=beadPost,line width=2.4pt,line cap=round},
  gbknot/.style={circle,fill=beadPost,inner sep=1.5pt},
  gbbead/.style={circle,draw=beadRim,line width=.5pt,minimum size=5.25mm,
    inner sep=0pt,font=\scriptsize,text=solarInk},
  gbpoly/.style={draw=beadRim,line width=.5pt,inner sep=0pt,minimum size=5.6mm},
  gbspacer/.style={rectangle,draw=ArtRed,fill=ArtRed,minimum width=1.3mm,
    minimum height=4.4mm,inner sep=0pt},
  gbspacer thin/.style={rectangle,draw=ArtRed!65,fill=ArtRed!65,
    minimum width=.6mm,minimum height=3.4mm,inner sep=0pt},
}
% Coordinates are evaluated first, so a caller may pass any pgfmath expression
% (a calc coordinate cannot nest the parentheses of cos(...)).
% \gbead{x}{y}{fill}{label}: one bead with a glass highlight.
\newcommand{\gbead}[4]{%
  \pgfmathsetmacro{\gbx}{#1}\pgfmathsetmacro{\gby}{#2}%
  \node[gbbead,fill=#3] at (\gbx,\gby) {#4};
  \fill[white,opacity=.55] ($(\gbx,\gby)+(-.075cm,.085cm)$) circle (1.35pt);}
% \gbeadw{x}{y}{fill}{label}: the same with a white numeral, for dark beads.
\newcommand{\gbeadw}[4]{%
  \pgfmathsetmacro{\gbx}{#1}\pgfmathsetmacro{\gby}{#2}%
  \node[gbbead,fill=#3,text=white] at (\gbx,\gby) {#4};
  \fill[white,opacity=.55] ($(\gbx,\gby)+(-.075cm,.085cm)$) circle (1.35pt);}
% \gbeadrim{x}{y}{rim colour}{fill}{label}: a bead whose rim names a channel.
\newcommand{\gbeadrim}[5]{%
  \pgfmathsetmacro{\gbx}{#1}\pgfmathsetmacro{\gby}{#2}%
  \node[gbbead,fill=#4,draw=#3,line width=.9pt] at (\gbx,\gby) {#5};
  \fill[white,opacity=.55] ($(\gbx,\gby)+(-.075cm,.085cm)$) circle (1.35pt);}
% \gbugle{x1}{x2}{y}{fill}{text colour}{label}: one bead held from the
% position centred at x1 to the position centred at x2.
\newcommand{\gbugle}[6]{%
  \pgfmathsetmacro{\gbxa}{#1}\pgfmathsetmacro{\gbxb}{#2}\pgfmathsetmacro{\gby}{#3}%
  \path[draw=beadRim,line width=.5pt,fill=#4,rounded corners=2.3mm]
    ($(\gbxa,\gby)+(-.2625cm,-.25cm)$) rectangle ($(\gbxb,\gby)+(.2625cm,.25cm)$);
  \fill[white,opacity=.42,rounded corners=.5mm]
    ($(\gbxa,\gby)+(-.13cm,.09cm)$) rectangle ($(\gbxb,\gby)+(.13cm,.165cm)$);
  \node[font=\scriptsize,text=#5] at ($(\gbxa,\gby)!.5!(\gbxb,\gby)$) {#6};}
% \nbead{x}{y}{n}{fill}: a numeral bead, n sides for the value n; a ring for
% 0, a disc for 1, a lens for 2.  In Z/8 the octagon is a full turn, 8 = 0.
\newcommand{\nbead}[4]{%
  \pgfmathsetmacro{\gbx}{#1}\pgfmathsetmacro{\gby}{#2}%
  \ifnum#3=0
    \draw[draw=beadRim,line width=.5pt,fill=#4,even odd rule]
      (\gbx,\gby) circle (.27cm) (\gbx,\gby) circle (.11cm);
    \fill[white,opacity=.5] ($(\gbx,\gby)+(-.07cm,.09cm)$) circle (1.1pt);
  \else\ifnum#3=1
    \node[gbbead,fill=#4] at (\gbx,\gby) {};
    \fill[white,opacity=.55] ($(\gbx,\gby)+(-.075cm,.085cm)$) circle (1.35pt);
  \else\ifnum#3=2
    \path[draw=beadRim,line width=.5pt,fill=#4]
      ($(\gbx,\gby)+(-.25cm,0)$) .. controls ($(\gbx,\gby)+(-.12cm,.22cm)$)
      and ($(\gbx,\gby)+(.12cm,.22cm)$) .. ($(\gbx,\gby)+(.25cm,0)$)
      .. controls ($(\gbx,\gby)+(.12cm,-.22cm)$)
      and ($(\gbx,\gby)+(-.12cm,-.22cm)$) .. cycle;
  \else
    \node[gbpoly,regular polygon,regular polygon sides=#3,fill=#4] at (\gbx,\gby) {};
  \fi\fi\fi}

% ------------------------------------------------- the pointer law (chapter 5)
% Added 2026-09-23 for Part A of chapter 5.  A pointer has three positions,
% numbered 0, 1, 2 anticlockwise from the page's up, a third of a turn apart;
% the hand of position k points at 90+120k degrees.  Corners of the running
% square are v_0..v_3 anticlockwise from the lower left.  Nothing drawn by
% these macros is data: no rule reads a hand.
% \lawdial{x}{y}{position}{rim colour}{radius}: one pointer.
\newcommand{\lawdial}[5]{%
  \begin{scope}[shift={(#1,#2)}]
    \path[draw=#4,thick,fill=white] (0,0) circle (#5);
    \foreach \k in {0,1,2}{\fill[ArtGray] ({90+120*\k}:#5) circle (.8pt);}
    \draw[ArtNavy,line width=1.1pt,-{Latex[length=1.3mm]}] (0,0)--({90+120*#3}:{.72*#5});
  \end{scope}}
% \lawsquare{x}{y}{side}{p0}{p1}{p2}{p3}{fill}{radius}: a state of the
% square, v_0 at (x,y).
\newcommand{\lawsquare}[9]{%
  \begin{scope}[shift={(#1,#2)}]
    \fill[#8] (0,0) rectangle (#3,#3);
    \draw[ArtNavy!45,thick] (0,0) rectangle (#3,#3);
    \lawdial{0}{0}{#4}{ArtNavy}{#9}\lawdial{#3}{0}{#5}{ArtNavy}{#9}
    \lawdial{#3}{#3}{#6}{ArtNavy}{#9}\lawdial{0}{#3}{#7}{ArtNavy}{#9}
  \end{scope}}
% \lawedges{x}{y}{side}: the oriented edges e_1..e_4 of the square, e_i from
% v_{i-1} to v_i, drawn clear of pointers of radius .24.
\newcommand{\lawedges}[3]{%
  \draw[ArtNavy!55,thick,-{Latex[length=1.8mm]}] (#1+.32,#2)--(#1+#3-.32,#2);
  \draw[ArtNavy!55,thick,-{Latex[length=1.8mm]}] (#1+#3,#2+.32)--(#1+#3,#2+#3-.32);
  \draw[ArtNavy!55,thick,-{Latex[length=1.8mm]}] (#1+#3-.32,#2+#3)--(#1+.32,#2+#3);
  \draw[ArtNavy!55,thick,-{Latex[length=1.8mm]}] (#1,#2+#3-.32)--(#1,#2+.32);}
% Numbers in place of pointer hands (author, 2026-09-27: "Replace pointer
% arrows in every diagrams by numbers").  A corner's number is its pointer's
% position, counted in thirds of a turn anticlockwise from the page's up.
% \lawnum{x}{y}{number}{size}: one number in a circle.
\newcommand{\lawnum}[4]{\node[circle,draw=ArtNavy,thick,fill=white,inner sep=0pt,
  minimum size=#4,font=\scriptsize] at (#1,#2) {$#3$};}
% \lawnumsquare{x}{y}{side}{p0}{p1}{p2}{p3}{fill}{size}: a state of the
% square, v_0 at (x,y), with numbers at its corners.
\newcommand{\lawnumsquare}[9]{%
  \begin{scope}[shift={(#1,#2)}]
    \fill[#8] (0,0) rectangle (#3,#3);
    \draw[ArtNavy!45,thick] (0,0) rectangle (#3,#3);
    \lawnum{0}{0}{#4}{#9}\lawnum{#3}{0}{#5}{#9}\lawnum{#3}{#3}{#6}{#9}\lawnum{0}{#3}{#7}{#9}
  \end{scope}}
% \walkstate{x}{y}{side}{k}: a state of the walker, a small square centred at
% (x,y) with its token at the corner v_k, k = 0..3.
\newcommand{\walkstate}[4]{%
  \begin{scope}[shift={(#1,#2)}]
    \fill[ArtLightBlue!40] (-#3/2,-#3/2) rectangle (#3/2,#3/2);
    \draw[ArtNavy!45] (-#3/2,-#3/2) rectangle (#3/2,#3/2);
    \ifcase#4 \fill[ArtNavy] (-#3/2,-#3/2) circle (.075);
    \or \fill[ArtNavy] (#3/2,-#3/2) circle (.075);
    \or \fill[ArtNavy] (#3/2,#3/2) circle (.075);
    \or \fill[ArtNavy] (-#3/2,#3/2) circle (.075);\fi
  \end{scope}}
% The spring law (chapter 5, A tick in steps).  \springedge{x}{y}{length}{angle}:
% a spring from (x,y), clear of numbers of radius .23 at its ends.
% \springstate{x}{y}{positions}{velocities}: a state, v_0 at (x,y), side 1.2,
% positions in circles at the corners, velocities in boxes beside them.
\tikzset{springvel/.style={draw=ArtOrange,thick,fill=ArtLightGold,rounded corners=1pt,
  inner sep=1pt,minimum size=3.6mm,font=\scriptsize}}
\newcommand{\springedge}[4]{%
  \begin{scope}[shift={(#1,#2)},rotate=#4]
    \draw[ArtNavy!55,thick] (.23,0)--(.34,0)
      \foreach \k in {1,...,6}{--({.34+(\k-.5)*(#3-.68)/6},{.07*(2*mod(\k,2)-1)})}
      --(#3-.34,0)--(#3-.23,0);
  \end{scope}}
\newcommand{\springstate}[4]{%
  \begin{scope}[shift={(#1,#2)}]
    \fill[ArtLightBlue!40] (0,0) rectangle (1.2,1.2);
    \springedge{0}{0}{1.2}{0}\springedge{1.2}{0}{1.2}{90}%
    \springedge{1.2}{1.2}{1.2}{180}\springedge{0}{1.2}{1.2}{270}%
    \foreach \v [count=\i from 0] in {#3}{%
      \pgfmathsetmacro{\cx}{{0,1.2,1.2,0}[\i]}\pgfmathsetmacro{\cy}{{0,0,1.2,1.2}[\i]}%
      \lawnum{\cx}{\cy}{\v}{4.4mm}}%
    \foreach \v [count=\i from 0] in {#4}{%
      \pgfmathsetmacro{\bx}{{-.38,1.58,1.58,-.38}[\i]}\pgfmathsetmacro{\by}{{-.3,-.3,1.5,1.5}[\i]}%
      \node[springvel] at (\bx,\by) {$\v$};}%
  \end{scope}}
% Two squares that share an edge (chapter 5, from section 5.5 on; 2026-09-28).
% Corners v_0=(0,0), v_1=(s,0), v_2=(s,s), v_3=(0,s) are the square of
% Part A; v_4=(2s,0) and v_5=(2s,s) are added on the right.  Part A's edges
% e_1..e_4 keep their direction, e_i from v_{i-1} to v_i, and the right
% square carries on anticlockwise: e_5 from v_1 to v_4, e_6 from v_4 to v_5,
% e_7 from v_5 to v_2.  Every outer edge points anticlockwise round the
% outside, and the shared edge e_2 points up.
% \tsqcoords{s}: name the six corners (v0)..(v5) at the scope's origin.
\newcommand{\tsqcoords}[1]{%
  \coordinate (v0) at (0,0); \coordinate (v1) at (#1,0); \coordinate (v2) at (#1,#1);
  \coordinate (v3) at (0,#1); \coordinate (v4) at (2*#1,0); \coordinate (v5) at (2*#1,#1);}
% \tsqfill{s}: the two faces, the left one blue and the right one gold, and the
% seven edges as thin lines.
\newcommand{\tsqfill}[1]{%
  \fill[ArtLightBlue!80] (0,0) rectangle (#1,#1);
  \fill[ArtLightGold] (#1,0) rectangle (2*#1,#1);
  \draw[ArtNavy!45,thick] (0,0) rectangle (2*#1,#1); \draw[ArtNavy!45,thick] (#1,0)--(#1,#1);}
% \tsqlink[style]{from}{to}{number}{placement}: an oriented link and its number;
% an empty placement puts the number on the link.
\newcommand{\tsqlink}[5][]{%
  \draw[ArtNavy,thick,-{Latex[length=1.9mm]},#1] ($(#2)!.2!(#3)$)--($(#2)!.8!(#3)$)
    node[midway,#5,plain label,inner sep=1pt]{#4};}
% \tsqnum{corner}{entry}{size} and \tsqsix{v_0..v_5 entries}{size}: entries at
% the corners.
\newcommand{\tsqnum}[3]{\node[circle,draw=ArtNavy,thick,fill=white,inner sep=0pt,
  minimum size=#3,font=\scriptsize] at (#1) {$#2$};}
\newcommand{\tsqsix}[7]{\tsqnum{v0}{#1}{#7}\tsqnum{v1}{#2}{#7}\tsqnum{v2}{#3}{#7}%
  \tsqnum{v3}{#4}{#7}\tsqnum{v4}{#5}{#7}\tsqnum{v5}{#6}{#7}}
\tikzset{tsqvtx/.style={circle,fill=ArtNavy,inner sep=1.4pt},
  turnL/.style={draw=ArtBlue,line width=1.1pt,-{Latex[length=1.7mm]}},
  turnR/.style={draw=ArtOrange,line width=1.1pt,-{Latex[length=1.7mm]}}}
% Runs strung as temporal tapestries (chapters 5 and 6; 2026-09-28, the author:
% "consitently use temporal tapestries with time runing left to right").  A run
% is a Hesse frame: one wire per cell, in the canvas's order from the top; one
% rung per tick, or per slot for a law in steps, time running left to right;
% posts at both ends.  A control row sits above the frame and the readings
% below it, as readouts and not wires.  A bar holds one value over consecutive
% rungs (PROSE-RULES section 9).  Nothing drawn by these macros is data.
% \tapframe[extra]{x0}{dx}{n}{ytop}{dy}{wires}: rungs 0..n at x0+dx*t, wires down
% from ytop; the optional extra moves the right post out, so that two contributions
% can stand side by side at the last rung.
\newcommand{\tapframe}[7][0]{%
  \pgfmathsetmacro{\tapxa}{#2-.45}\pgfmathsetmacro{\tapxb}{#2+#3*#4+.45+#1}%
  \pgfmathsetmacro{\tapya}{#5+.34}\pgfmathsetmacro{\tapyb}{#5-#6*(#7-1)-.34}%
  \foreach \t in {0,...,#4}{\draw[solarLine,line width=.3pt] ({#2+#3*\t},\tapyb)--({#2+#3*\t},\tapya);}
  \foreach \r in {1,...,#7}{\draw[gbwire] (\tapxa,{#5-#6*(\r-1)})--(\tapxb,{#5-#6*(\r-1)});}
  \draw[gbpost] (\tapxa,\tapyb)--(\tapxa,\tapya);
  \draw[gbpost] (\tapxb,\tapyb)--(\tapxb,\tapya);}
% \taprungs{x0}{dx}{n}{y}: the rung numerals.
\newcommand{\taprungs}[4]{\foreach \t in {0,...,#3}{%
  \node[font=\scriptsize,text=solarMuted] at ({#1+#2*\t},#4) {$\t$};}}
% \tapbead{x}{y}{v}: a number 0, 1 or 2 as a bead, white, teal or purple, the
% colours of fig:square-gauge-beads; \taprow{x0}{dx}{y}{values}: one wire's beads.
\newcommand{\tapbead}[3]{%
  \ifcase#3\relax \gbead{#1}{#2}{white}{$0$}\or \gbead{#1}{#2}{ArtTeal!42}{$1$}\or
  \gbead{#1}{#2}{ArtPurple!42}{$2$}\fi}
\newcommand{\taprow}[4]{\foreach \v [count=\t from 0] in {#4}{\tapbead{#1+#2*\t}{#3}{\v}}}
% Letters standing for numbers, one colour each, never a residue's colour; b,
% the letter the two-squares law keeps at v_1, is gold.
\colorlet{tapLa}{ArtBlue!38}\colorlet{tapLb}{ArtGold!70}\colorlet{tapLc}{ArtPurple!35}
\colorlet{tapLd}{ArtRed!38}\colorlet{tapLe}{ArtGreen!38}\colorlet{tapLf}{ArtOrange!42}
% \tapletter{x}{y}{letter} and \tapletters{x0}{dx}{y}{letters}.
\newcommand{\tapletter}[3]{\gbead{#1}{#2}{tapL#3}{$#3$}}
\newcommand{\tapletters}[4]{\foreach \v [count=\t from 0] in {#4}{\tapletter{#1+#2*\t}{#3}{\v}}}
% \tapread{x1}{x2}{y}{fill}{label}: a reading held from x1 to x2, gold rim;
% \tapreadbead{x}{y}{fill}{label}: a reading at one rung.
\newcommand{\tapread}[5]{%
  \pgfmathsetmacro{\gbxa}{#1}\pgfmathsetmacro{\gbxb}{#2}\pgfmathsetmacro{\gby}{#3}%
  \path[draw=ArtGold!80!black,line width=.9pt,fill=#4,rounded corners=2.3mm]
    ($(\gbxa,\gby)+(-.2625cm,-.25cm)$) rectangle ($(\gbxb,\gby)+(.2625cm,.25cm)$);
  \fill[white,opacity=.42,rounded corners=.5mm]
    ($(\gbxa,\gby)+(-.13cm,.09cm)$) rectangle ($(\gbxb,\gby)+(.13cm,.165cm)$);
  \node[font=\scriptsize,text=ArtNavy] at ($(\gbxa,\gby)!.5!(\gbxb,\gby)$) {#5};}
\newcommand{\tapreadbead}[4]{\gbeadrim{#1}{#2}{ArtGold!80!black}{#3}{#4}}
% \tapheld{x1}{x2}{y}{channel colour}{residue}: a channel's value held from x1
% to x2 (a round bead when x1 = x2), rim in the channel's colour.
\newcommand{\tapheld}[5]{%
  \pgfmathsetmacro{\gbxa}{#1}\pgfmathsetmacro{\gbxb}{#2}\pgfmathsetmacro{\gby}{#3}%
  \path[draw=#4,line width=.9pt,fill=#4!30,rounded corners=2.3mm]
    ($(\gbxa,\gby)+(-.2625cm,-.25cm)$) rectangle ($(\gbxb,\gby)+(.2625cm,.25cm)$);
  \fill[white,opacity=.42,rounded corners=.5mm]
    ($(\gbxa,\gby)+(-.13cm,.09cm)$) rectangle ($(\gbxb,\gby)+(.13cm,.165cm)$);
  \node[font=\scriptsize,text=solarInk] at ($(\gbxa,\gby)!.5!(\gbxb,\gby)$) {$#5$};}
% \tapsmall{x}{y}{fill}{label}: a small bead, for the prism.
\newcommand{\tapsmall}[4]{\node[gbbead,minimum size=3.7mm,font=\tiny,fill=#3] at ({#1},{#2}) {#4};}
% \tapstep{x}{y}{L or R}: a bead of the two squares' score; \tapphase{x}{y}{K or M}: of
% the Maxwell score.
\newcommand{\tapstep}[3]{\if#3L\node[draw=ArtBlue,fill=ArtBlue!16,rounded corners=2pt,
    font=\scriptsize\bfseries,text=ArtBlue,minimum width=4.6mm,inner sep=1pt] at ({#1},{#2}) {L};%
  \else\node[draw=ArtOrange,fill=ArtOrange!20,rounded corners=2pt,font=\scriptsize\bfseries,
    text=ArtOrange!85!black,minimum width=4.6mm,inner sep=1pt] at ({#1},{#2}) {R};\fi}
\newcommand{\tapphase}[3]{\if#3K\node[draw=ArtOrange,fill=ArtOrange!20,rounded corners=2pt,
    font=\scriptsize\bfseries,text=ArtOrange!85!black,minimum width=4.6mm,inner sep=1pt] at ({#1},{#2}) {K};%
  \else\node[draw=ArtTeal,fill=ArtTeal!20,rounded corners=2pt,font=\scriptsize\bfseries,
    text=ArtTeal,minimum width=4.6mm,inner sep=1pt] at ({#1},{#2}) {M};\fi}
\tikzset{
  tapthread/.style={line width=2.2pt,line cap=round,opacity=.75},
  tapwindow/.style={draw=ArtOrange,line width=1.2pt,rounded corners=3.2mm,inner sep=0pt},
  tapnote/.style={font=\scriptsize,align=left,anchor=north west,inner sep=1pt},
  taphead/.style={font=\scriptsize\bfseries,anchor=west,text=ArtNavy}}

% Coherent runs and scores strung as tapestries (Chapters 7-11; 2026-09-28, the
% author: "audit chapter 7 ... for dynamics that should be temporal tapestries";
% "go with your recommendations, implement and publish").  A coherent run is strung
% over its complete labels: one wire per label, the cells of a coefficient table.  A
% bead is one history's contribution, coloured and numbered by its phase mark
% (ph0-ph7); its size [1,g] is printed once below the rung; a grey thread is one
% history.  A score's footprint is drawn as cells over the beads, dashed where a slot
% reads the cell and shaded where it writes it; gold marks a checkpoint rung.
% Nothing drawn by these macros is data.
\tikzset{
  tapgate/.style={draw=ArtNavy!65,fill=ArtNavy!7,rounded corners=2pt,font=\scriptsize,
    inner sep=1.4pt,minimum height=4.3mm,text=ArtNavy},
  tapcell read/.style={draw=ArtNavy!75,dashed,line width=.6pt,rounded corners=1.5pt},
  tapcell K/.style={fill=ArtOrange,fill opacity=.32,draw=ArtOrange!80!black,line width=.4pt,
    rounded corners=1.5pt},
  tapcell M/.style={fill=ArtTeal,fill opacity=.30,draw=ArtTeal!80!black,line width=.4pt,
    rounded corners=1.5pt},
  tapcell W/.style={fill=ArtNavy,fill opacity=.18,draw=ArtNavy!60,line width=.4pt,
    rounded corners=1.5pt},
  tapcheck/.style={draw=ArtGold!85!black,line width=1.3pt},
  taphistory/.style={draw=ArtGray!55,line width=1.6pt,line cap=round,opacity=.75},
  tapquark/.style={draw=ArtGold,line width=2.2pt,line cap=round,opacity=.85},
  tapbrace/.style={decorate,decoration={brace,amplitude=4pt},draw=ArtGray,thick},
  taplabel/.style={font=\scriptsize,anchor=east,text=ArtNavy},
  taptitle/.style={font=\small\bfseries,text=ArtNavy,anchor=west}}
% \tapcell{x0}{dx}{slot}{y}{style}: one wire's cell in one slot, drawn over the beads.
\newcommand{\tapcell}[5]{\path[#5] ({#1+#2*(#3)+.06},{#4-.3}) rectangle ({#1+#2*((#3)+1)-.06},{#4+.3});}
% \tapcheck{x}{ytop}{ybottom}: a checkpoint rung; \tapchecktick: its two ends only,
% for a frame whose bars would hide the rung.
\newcommand{\tapcheck}[3]{\draw[tapcheck] ({#1},{#2}) -- ({#1},{#3});}
\newcommand{\tapchecktick}[3]{\draw[tapcheck] ({#1},{#2}) -- ({#1},{#2+.2}) ({#1},{#3}) -- ({#1},{#3-.2});}
% \tapfluxgrid{x}{y}{s or r}{b_L}{b_R}: the sixteen flux labels of two faces as a 4x4
% glyph, top left at (x,y): all dark where the field is spread (s), one dark dot at
% (b_L,b_R) where it is a single label (r).
\newcommand{\tapfluxgrid}[5]{%
  \foreach \i in {0,...,3}\foreach \j in {0,...,3}{%
    \fill[ArtGray!45] ({#1+.11*\j},{#2-.11*\i}) circle (.9pt);}%
  \if#3s\foreach \i in {0,...,3}\foreach \j in {0,...,3}{%
    \fill[ArtNavy!70] ({#1+.11*\j},{#2-.11*\i}) circle (1.25pt);}\fi
  \if#3r\fill[ArtNavy] ({#1+.11*#5},{#2-.11*#4}) circle (1.8pt);\fi}



\captionnamefont{\small\bfseries\color{ArtNavy}}
\captiontitlefont{\small\artfbkfloat}
\setlength{\abovecaptionskip}{5pt}
\setlength{\belowcaptionskip}{2pt}

\crefname{enumi}{exercise}{exercises}
\Crefname{enumi}{Exercise}{Exercises}
\crefname{theorem}{theorem}{theorems}
\crefname{proposition}{proposition}{propositions}
\crefname{lemma}{lemma}{lemmas}
\crefname{corollary}{corollary}{corollaries}
\crefname{definition}{definition}{definitions}
\crefname{example}{example}{examples}
\crefname{principle}{principle}{principles}

\newtcolorbox{literatureconnection}[1][]{
  enhanced,breakable,colback=ArtPaper,colframe=ArtTeal,
  boxrule=0.5pt,arc=2pt,left=7pt,right=7pt,top=6pt,bottom=6pt,
  title={Connections and sources},fonttitle=\bfseries,
  fontupper=\small,artfbk=literatureconnection,#1}
% The optional categorical track (2026-09-22, planning/category-cafe-optional-
% 2026-09-22). A section, exercise or box marked \optionalmark is optional: it
% reads the book's constructions in the language of category theory, and no
% unmarked text depends on it. tools/verify_reading_map.py masks the box the
% way it masks exercises. The box enters use with the later packets.
\newcommand*{\optionalmark}{\ensuremath{\lozenge}}
\newtcolorbox{categoricalview}[1][]{
  enhanced,breakable,colback=ArtPaper,colframe=ArtPurple,
  boxrule=0.5pt,arc=2pt,left=7pt,right=7pt,top=6pt,bottom=6pt,
  title={Categorical view \optionalmark},fonttitle=\bfseries,
  fontupper=\small,artfbk=categoricalview,#1}
% The optional type-theoretic track (2026-09-28, planning/type-theory-view-
% 2026-09-28): Section 1.11 states Martin-Lof's rules behind the book's notation,
% and one box in each of Sections 2.5, 5.6, 6.4 and 7.4 reads that chapter's
% construction in them. Marked like the categorical track; nothing unmarked
% cites it, and tools/verify_reading_map.py masks the box the same way.
\newtcolorbox{typetheoryview}[1][]{
  enhanced,breakable,colback=ArtPaper,colframe=ArtNavy,
  boxrule=0.5pt,arc=2pt,left=7pt,right=7pt,top=6pt,bottom=6pt,
  title={Type-theoretic view \optionalmark},fonttitle=\bfseries,
  fontupper=\small,before upper={\setlength{\parskip}{4pt}},artfbk=typetheoryview,#1}
\newcommand{\Id}{\operatorname{Id}}
\newcommand{\refl}{\mathsf{refl}}
\newcommand{\Hist}{\mathsf{Hist}}
\newcommand{\App}{\mathsf{App}}
\newcommand{\coef}{\operatorname{coef}}
\newcommand{\Contact}{\mathsf{Contact}}
\newcommand{\NoRow}{\mathsf{NoRow}}
\newcommand{\stepto}[1]{\xrightarrow{\;#1\;}}
\newcommand{\rulename}[1]{\textsc{#1}}

% ------------------------------------------------------------------ float breaks
% The author, 2026-09-29: "the figure never break a plate" and "the figure never
% breaks a theorem". A breakable box, or a theorem-like statement, that runs on
% to a later page must go on there with no float between its parts: no float page
% between them, no float below its first part, and none above its last. Nothing
% in the source says where a float will land, so positions are recorded at
% shipout, in \jobname.fbk, and tools/find_float_breaks.py reads them.
% Every mark sits in horizontal mode, inside a line of text: a mark on a vertical
% list would let \parskip in ahead of the next paragraph and hide \lastskip from
% \addvspace, and the pages would move. So a statement's start is marked in its
% head and a box's in its title; the end is marked on the last line, or, when the
% statement ends in a list, on the list's last line (an L record); a float is
% marked on its caption. An end mark first removes the space before it, as \par
% would have: a space left in front of a mark can push a word onto a new line.
% Fields: record, kind, serial, page count, folio, height from the foot of the
% page in sp, name. Defining \artfbknoT, \artfbknoE, \artfbknoL, \artfbknoF or
% \artfbknoB before this file switches off the head, end, list-end, float or
% box marks, for finding a mark that moves the pages (planning/simpler-models-
% 2026-09-29/harness): with all five on, the pages are those without any.
\newwrite\artfbk
\AtBeginDocument{\immediate\openout\artfbk=\jobname.fbk\relax}
\newcounter{artfbk}
\def\artfbkid{0}
\newcommand\artfbknew{\stepcounter{artfbk}\edef\artfbkid{\arabic{artfbk}}}
\newcommand\artfbkmark[3]{%
  \pdfsavepos
  \begingroup\edef\x{\endgroup\write\artfbk{#1 #2 \artfbkid\space
    \noexpand\the\noexpand\ReadonlyShipoutCounter\space\noexpand\thepage\space
    \noexpand\the\noexpand\pdflastypos\space #3}}\x}
\makeatletter
\let\artfbkpending\@empty
\newcommand\artfbkhead{\ifdefined\artfbknoT\else\ifx\artfbkpending\@empty\else
  \global\let\artfbkpending\@empty\artfbkmark{S}{thm}{\@currentlabel}\fi\fi}
\def\artfbkthm{\artfbknew\gdef\artfbkpending{x}}
\def\artfbkend#1{\ifdefined\artfbknoE\else\ifhmode\ifnum\lastnodetype=-1 \else
  \unskip\artfbkmark{E}{#1}{}\fi\fi\fi}
\@for\artfbkenv:={theorem,proposition,lemma,corollary,definition,example,principle}\do{%
  \expandafter\AtBeginEnvironment\expandafter{\artfbkenv}{\artfbkthm}%
  \expandafter\AtEndEnvironment\expandafter{\artfbkenv}{\artfbkend{thm}}}
\@for\artfbkenv:={enumerate,itemize,description,center,quote,quotation,flushleft,flushright}\do{%
  \expandafter\AtEndEnvironment\expandafter{\artfbkenv}{\ifdefined\artfbknoL\else\ifhmode\unskip\artfbkmark{L}{any}{}\fi\fi}}
% a statement or box that ends in a display leaves an empty paragraph after it,
% which TeX discards; a mark there would print it as a blank line, so the end is
% marked inside the display instead
\@for\artfbkenv:={equation,equation*,align,align*,gather,gather*,multline,multline*,mathpar}\do{%
  \expandafter\AtEndEnvironment\expandafter{\artfbkenv}{\ifdefined\artfbknoL\else\artfbkmark{L}{any}{}\fi}}
\def\artfbk@figure{figure}\def\artfbk@table{table}
\newcommand\artfbkfloat{\ifdefined\artfbknoF\else\artfbkfloatmark\fi}
\newcommand\artfbkfloatmark{\ifnum0\ifx\@currenvir\artfbk@figure1\fi\ifx\@currenvir\artfbk@table1\fi>0
  \begingroup\def\artfbkid{0}\artfbkmark{F}{float}{\@currentlabel}\endgroup\fi}
\makeatother
\tcbset{artfbk/.style={before title app={\ifdefined\artfbknoB\else\leavevmode\artfbkmark{S}{box}{#1}\fi},
  after upper pre={\artfbkend{box}}}}
\makeatletter\@for\artfbkenv:={picturefirst,worked,formalcorner,design,chaptergoals,literatureconnection,categoricalview,typetheoryview}\do{%
  \expandafter\AtBeginEnvironment\expandafter{\artfbkenv}{\artfbknew}}
\makeatother

